\subsection{Discusión de los resultados}

Los resultados obtenidos demuestran que el sistema LPR-Sentinel es funcional y alcanza un rendimiento competitivo en condiciones reales, especialmente considerando las limitaciones de hardware de la plataforma objetivo (Raspberry Pi 4B). La tasa de éxito integral del 77.5% sobre imágenes reales capturadas en condiciones variadas indica que el sistema es suficientemente robusto para muchas aplicaciones prácticas, como el control de acceso a estacionamientos residenciales (donde las condiciones de captura suelen ser controladas) o el monitoreo de flotas empresariales en horario diurno. La latencia media de 121 ms permite un procesamiento casi en tiempo real, y la modularidad del sistema facilita la actualización de componentes individuales sin necesidad de rediseñar el pipeline completo.

Las principales limitaciones identificadas son la sensibilidad a rotaciones extremas (superiores a 60°), a desenfoques de movimiento severos y a condiciones de iluminación extremadamente desfavorables (sombras profundas o contraluz intensa). Estas limitaciones sugieren líneas de trabajo futuro, como la inclusión de transformaciones de aumentación más agresivas durante el entrenamiento (rotaciones de hasta 90°, desenfoques direccionales de mayor intensidad) o la incorporación de un módulo de preprocesamiento adaptativo que normalice la iluminación antes de la inferencia. No obstante, el sistema alcanza su objetivo principal: demostrar la viabilidad de un pipeline completo de reconocimiento de matrículas mexicanas basado en datos sintéticos, con componentes desplegables en formato ONNX sobre hardware de bajo costo.